% !TeX root = ../../paper.tex

\subsection{Interpretation}
\label{sec:discussion-findings}

Taken at face value, the results are a strong endorsement of passive haptics: every participant who played both versions reported more presence with the proxies, on every subscale, and the differences are far larger than what presence manipulations typically produce. The direction is consistent with the classic findings~\cite{Hoffman.1998, Insko.2001} and extends them to deliberately aversive materials---confirming a dreaded touch strengthened the illusion rather than breaking it. That the largest item-level gains sit exactly where a haptic proxy should act (feeling present in the space, consistency with real-world experience, and reduced awareness of the real surroundings) adds internal plausibility: the effect concentrates on the items the manipulation targets, and stays smallest where no proxy could help (a virtual cabin seeming \emph{more} real than reality).

The size of the effect still deserves suspicion rather than celebration. Participants could not be blinded---one condition involves visibly laid-out props and a researcher with a spray bottle---and in a within-subject design this invites demand characteristics: participants may have rated the contrast between conditions rather than each experience on its own. The strikingly low baseline scores point the same way. Subscale means of 2--3 are low for a functioning VR experience, which suggests participants used the questionnaire comparatively, pushing the baseline down after (or in anticipation of) the richer haptic version. A between-subject replication, in which each participant sees only one condition, would separate the genuine haptic effect from this contrast effect; we would expect the difference to shrink, though the concentration of gains on touch-related items gives some confidence it would not vanish.

We also note what the data cannot show. The IPQ measures presence, not fear or disgust, so whether the proxies made the experience \emph{scarier}---our original motivation---remains formally unmeasured. The involvement subscale's gain of $+2.65$ is at least suggestive: participants with proxies reported being far less aware of the room, the researchers, and the sounds around them, which in a horror context is where dread lives.
